\documentclass[11pt]{article}

\usepackage[T1]{fontenc}
\usepackage{lmodern}
\usepackage[a4paper,margin=27mm]{geometry}
\usepackage{amsmath,amssymb,amsthm,mathtools,mathrsfs}
\usepackage{microtype}
\usepackage{enumitem}
\usepackage{xcolor}
\usepackage[colorlinks=true,linkcolor=blue!45!black,citecolor=blue!45!black,
  urlcolor=blue!55!black]{hyperref}
\setlength{\emergencystretch}{2em}

\newtheorem{theorem}{Theorem}[section]
\newtheorem{proposition}[theorem]{Proposition}
\newtheorem{lemma}[theorem]{Lemma}
\newtheorem{corollary}[theorem]{Corollary}
\newtheorem*{maintheorem}{Prime-valuation realization theorem}
\theoremstyle{definition}
\newtheorem{definition}[theorem]{Definition}
\theoremstyle{remark}
\newtheorem{remark}[theorem]{Remark}

\newcommand{\C}{\mathbb C}
\newcommand{\R}{\mathbb R}
\newcommand{\N}{\mathbb N}
\newcommand{\cF}{\mathcal F}
\newcommand{\cH}{\mathcal H}
\newcommand{\cP}{\mathcal P}
\newcommand{\cD}{\mathcal D}
\newcommand{\cR}{\mathcal R}
\newcommand{\cZ}{\mathcal Z}
\newcommand{\ord}{\operatorname{ord}}
\newcommand{\supp}{\operatorname{supp}}
\newcommand{\Span}{\operatorname{span}}
\newcommand{\Thetaas}{\mathop{\Theta}}

\title{A Universal Asymptotic-Type Superselection Law\\
\large Prime-Valuation Reflection and the Riemann Boundary}
\author{Michael Haufschild}
\date{}

\begin{document}

\maketitle

\begin{abstract}
Let \(a_N\to0\).  An all-orders-flat germ \(x_N=o(a_N^m)\) for every \(m\)
cannot be sewn to a nonzero finite-grade germ \(y_N=\Theta(a_N^d)\) by any bounded
operator family of fixed finite Rees order, in any dimension.  More sharply, every exact
sewing \(R_Nx_N=y_N\) satisfies
\[
 a_N^{-K}=o(\|R_N\|)\qquad\text{for every }K.
\]
This asymptotic-type superselection law is independent of arithmetic, rank, and choice of
normed spaces.

The prime-valuation Euler source gives a canonical arithmetic realization in which reflection
about the Riemann critical line creates the two required types before any zero is selected.  Let
\[
 \mu_{p,\varepsilon}(k)=
 (1-p^{-1+2\varepsilon})p^{-k(1-2\varepsilon)}
 \qquad(k\geq0)
\]
be the geometric prime-valuation law, and let \(\mathcal R_N(c)\) be the Hellinger affinity
between the reflected product states
\(\bigotimes_{p\leq N}\mu_{p,c/\log N}\) and
\(\bigotimes_{p\leq N}\mu_{p,-c/\log N}\).  We prove
\[
 \mathcal R_N(c)^2\longrightarrow e^{-\Phi(c)},\qquad
 \Phi(c)=4\int_0^1\frac{\sinh^2(cu)}u\,du .
\]
The local squared factor is exactly
\[
 1-\frac{4p^{-1}\sinh^2(c\log p/\log N)}{(1-p^{-1})^2}.
\]
The same \(\Phi\) is the limit of the symmetric prime-log Poisson/Wick cumulant
\[
 K_N(2c/\log N)=
 \sum_{p\leq N}\frac{2\{\cosh(2c\log p/\log N)-1\}}p ,
\]
and in fact \(\mathcal R_N(c)^2e^{K_N(2c/\log N)}\to1\).  Thus one nonlinear
profile is simultaneously the Hellinger action of reflected valuation states and the
L\'evy--Wick cumulant of their prime clock.

The critical window has asymmetric outer faces.  Comparing a fixed displaced state with the
critical state gives the exact Euler-partition curvature
\[
 A_N(c)^2=\frac{Z_N(1-c)^2}{Z_N(1)Z_N(1-2c)},\qquad
 Z_N(\beta)=\prod_{p\leq N}(1-p^{-\beta})^{-1}.
\]
With \(a_N=Z_N(1)^{-1}\), one face has
\(A_N(c)^2/a_N\to\zeta(1-c)^2/\zeta(1-2c)\), whereas its reflection is flat
beyond every power of \(a_N\).  These are not two finite grades: they are different asymptotic
species.  Kakutani singularity collapses them; their Rees type recovers the lost orientation.

The profile \(\Phi\) is even, so its scalar action cannot itself carry side.  We prove the
complementary signed law: a rank-two pointed Fisher connection is odd under tangent reversal,
and its canonical cutoff Ward trace equals its pairing with the double-cutoff curvature.  This
identifies an exact pre-divisor carrier of orientation rather than squaring it away.

The abstract theorem now applies: no finite-Rees-moderate operator can sew an off-line reflected
pair, and every exact sewing has an essential scale singularity.  Standard Weyl reflection
escapes only by modular retuning; divisor differentiation escapes only through a singular graded
correspondence.  For the completed zeta current, finite-aperture transport is already four-unit
nonexpansive, while the RH-complete condition is exactly the corresponding zero-aperture
Ward-tangent bound.  These unconditional results do not prove the Riemann hypothesis.  They
identify a previously missing pre-divisor invariant and reduce the viable source-internal route
to transport of this valuation type through a singular boundary tangent.
\end{abstract}

\medskip
\noindent\textbf{2020 Mathematics Subject Classification.}
Primary 11M26, 60E15; Secondary 11M41, 46B99, 12H05.

\noindent\textbf{Keywords.}
Riemann zeta function; prime valuations; Hellinger affinity; Kakutani dichotomy;
Rees filtration; Weyl intertwiner; modular drift; singular transport; hypertranscendence.

\section{Introduction}\label{sec:introduction}

The completed Riemann function
\[
 \Xi(s)=\frac12s(s-1)\pi^{-s/2}\Gamma(s/2)\zeta(s)
\]
satisfies
\[
 \Xi(s)=\Xi(1-s).
 \tag{1.1}\label{eq:functional-equation}
\]
At the level of zeros this pairs
\(\rho=\tfrac12+c+i\gamma\) with \(1-\rho=\tfrac12-c-i\gamma\).  The pairing is exact,
but its scalar statement is orientation-free: after the pair has been formed, it does not say
which member came from the right face \(c>0\) and which from the left face \(c<0\).

Many spectral and positivity approaches attempt to recover this missing sign after passing to a
Hilbert space, a trace, a determinant, a positive kernel, or a divisor count.  The present paper
starts one stage earlier.  An Euler product is first a product of local valuation laws.  We ask
whether reflection already changes the \emph{asymptotic type} of those laws before they are
assembled into a scalar meromorphic function.

The answer is unexpectedly rigid and has two scales.  For a prime \(p\) and a real radial
displacement \(|\varepsilon|<\tfrac12\), consider the geometric law
\[
 \mu_{p,\varepsilon}(k)=(1-r_p(\varepsilon))r_p(\varepsilon)^k,
 \qquad r_p(\varepsilon)=p^{-(1-2\varepsilon)}.
 \tag{1.2}\label{eq:valuation-law-intro}
\]
Reflection sends \(\varepsilon\) to \(-\varepsilon\).  At the shrinking displacement
\(\varepsilon_N=c/\log N\), let
\[
 \mathcal R_N(c)=\prod_{\substack{p\leq N\\p\ {\rm prime}}}
 \sum_{k\geq0}
 \sqrt{\mu_{p,c/\log N}(k)\mu_{p,-c/\log N}(k)}
 \tag{1.3}\label{eq:hellinger-intro}
\]
be the reflected Hellinger affinity.  The central theorem is the nonlinear crossover
\[
 \boxed{\quad
 \mathcal R_N(c)^2\longrightarrow e^{-\Phi(c)},\qquad
 \Phi(c)=4\int_0^1\frac{\sinh^2(cu)}u\,du .
 \quad}
 \tag{1.4}\label{eq:critical-profile-intro}
\]
The local squared factor is \(1-x_{p,N}(c)\), where
\[
 x_{p,N}(c)=
 \frac{4p^{-1}\sinh^2(c\log p/\log N)}{(1-p^{-1})^2}.
 \tag{1.5}\label{eq:local-defect-intro}
\]
The denominator is asymptotically invisible in aggregate, while Mertens's first theorem sends
the numerator sum to \(\Phi(c)\).  Uniformly \(x_{p,N}(c)\to0\), so logarithmic
tensorization turns addition into multiplication.

This profile already exists elsewhere in the source.  The symmetric prime-log cumulant
\[
 K_N(z)=\sum_{p\leq N}\frac{2\{\cosh(z\log p)-1\}}p
 \tag{1.6}\label{eq:wick-cumulant-intro}
\]
satisfies \(K_N(2c/\log N)\to\Phi(c)\), hence
\[
 \mathcal R_N(c)^2e^{K_N(2c/\log N)}\longrightarrow1.
 \tag{1.7}\label{eq:hellinger-wick-intro}
\]
This Hellinger--Wick reciprocity is the arithmetic finite-size law of the paper: information distance
between the reflected valuation sectors and fluctuation of the prime clock are two
representations of one action.

Here and below, ``universal'' refers first to the asymptotic-type theorem, which holds for
arbitrary asymptotic germs and real normed source and target spaces, and second to its cutoff-stable
realization inside the rational-prime valuation source.  It is not a claim that the displayed
prime profile holds for arbitrary Euler products or arbitrary reflection mechanisms.

The outer law is oriented.  Compare a fixed displaced state with the critical state:
\[
 A_N(c)=\prod_{p\leq N}
 \sum_{k\geq0}\sqrt{\mu_{p,0}(k)\mu_{p,c}(k)},\qquad
 a_N=\prod_{p\leq N}(1-p^{-1}).
 \tag{1.8}\label{eq:fixed-sector-intro}
\]
Then reflection does not merely change a finite exponent; it crosses a type boundary:
\[
 \boxed{\quad
 c<0:\ A_N(c)\asymp a_N^{1/2},\qquad
 c=0:\ A_N(c)=1,\qquad
 c>0:\ A_N(c)=o(a_N^m)\ \ \forall m.
 \quad}
 \tag{1.9}\label{eq:three-phase-intro}
\]
Writing \(Z_N(\beta)=\prod_{p\leq N}(1-p^{-\beta})^{-1}\), the exact curvature identity
\[
 A_N(c)^2=\frac{Z_N(1-c)^2}{Z_N(1)Z_N(1-2c)}
 \tag{1.10}\label{eq:partition-curvature-intro}
\]
gives \(A_N(c)^2/a_N\to\zeta(1-c)^2/\zeta(1-2c)\) for \(c<0\).
Thus every off-line reflected pair consists of one finite half-normal sector and one sector
flat beyond every finite Rees order.

There are four reasons this is more than another equivalent reformulation of the Riemann
hypothesis.

First, the critical-window theorem identifies an object that exists before the zeta divisor:
the same nonlinear action controls both information loss under reflection and fluctuation of an
infinitely divisible prime-log field.  Neither side is defined through zeros, and their
asymptotic reciprocity is an equality of source laws, not an RH criterion.

Second, it refines a classical theorem rather than decorating a divisor.  Kakutani's product
measure theorem~\cite{Kakutani1948} says that the infinite product laws for \(c=0\) and
\(c\neq0\) are mutually singular precisely when the Hellinger product vanishes.  In the present
family the product vanishes on \emph{both} off-critical faces.  The ordinary dichotomy therefore
collapses the two answers to the same symbol.  Formula \eqref{eq:three-phase-intro} recovers the
discarded orientation by resolving the rate at which singularity is reached.

Third, asymptotic type obeys an arithmetic-independent superselection law.  Suppose
\(a_N\to0\), \(x_N\) is flat beyond every power of \(a_N\),
\(\|y_N\|=\Theta(a_N^d)\) is nonzero of finite grade, and \(R_N:E\to F\) is a
family of bounded linear maps between arbitrary real normed spaces satisfying
\[
 R_Nx_N=y_N.
 \tag{1.11}\label{eq:moderate-intro}
\]
Then
\[
 a_N^{-K}=o(\|R_N\|)\qquad\text{for every }K.
 \tag{1.12}\label{eq:operator-invariance-intro}
\]
In particular, every operator family with one fixed finite Rees bound preserves the
all-orders-flat radical and cannot map it to a finite-grade germ.  A larger matrix, an
infinite-rank Hilbert space, or a bounded dilation does not alter the conclusion.  The
obstruction is not a missing finite renormalization but an essential scale singularity.

In the prime realization, if
\[
 S_N(c)=\frac{H_N(-c)}{H_N(c)}
\]
is the scalar sewing multiplier for \(0<c<\tfrac12\), then
\[
 a_N^{-K}=o\!\left(S_N(c)\right)
 \qquad\text{for every fixed }K.
\]
Every exact bounded linear sewing family has operator norm at least this costly in asymptotic
type.

Fourth, the theorem identifies a positive escape from that obstruction.  The rank-one rational
Weyl block has complementary eigenvalue
\[
 \lambda_p(z)=\frac{pz^2-1}{p-z^2}.
 \tag{1.13}\label{eq:weyl-intro}
\]
Its modulus defect factors exactly through \(|z|^4-1\).  It is unitary only on the unit circle;
away from it, reflection remains possible by changing modular class.  The no-sewing theorem
therefore separates \emph{internal type preservation} from \emph{external modular retuning}; it
does not mistake a regularity class for all possible reflection.

The universal action is even, and this exposes rather than solves the orientation problem at the
critical tangent.  A second exact law supplies the missing polarization.  The rank-two pointed
Fisher connection \(F_{\psi,\chi}\) changes sign with the source tangent \(\chi\), and its
cutoff Ward trace satisfies
\[
 \operatorname{Tr}\bigl([F_{\psi,\chi},e]^*[T,e]\bigr)
 =
 \operatorname{Tr}\bigl(F_{\psi,\chi}[[T,e],e^*]\bigr).
 \tag{1.14}\label{eq:signed-transgression-intro}
\]
Thus the prime action determines size, while an odd Fisher--Ward polarization can carry side
without waiting for a zero or squaring a normal coordinate.

Internal transport still forces a precise change of category.  The diagonal activation
\(\zeta(s)/\zeta(2s)\) has a local probability law whose nonzero event is squarefree occupancy
and whose conditional positive valuation is geometric.  At the critical exponent, occupancy is
prime-harmonic, all fixed tails of depth at least two are summable, and Fisher whitening returns
the fibre to prime-harmonic size.  Thus the archimedean Gamma gauge is not a missing bounded
source coordinate; it is the singular whitening of a rare event.

The same grade loss appears at a zero.  If a source overlap has local form \(B(t)=t^d\), then its
first numerator has form \(S(t)=dt^{d-1}\).  Positivity of
\[
 \begin{pmatrix}B(t)^2&S(t)\\S(t)&C(t)\end{pmatrix}
 \tag{1.15}\label{eq:gram-intro}
\]
forces \(C(t)\geq d^2/t^2\).  A bounded positive completion is impossible.  Division by the
overlap, often rejected as an unbounded defect, is exactly the operation required by the divisor
grade.

The paper closes the resulting circle.  The positive squarefree source supplies an explicit
M\"obius--Gamma graph: its unnormalized Fisher numerator divided by its own overlap reconstructs
the completed squarefree logarithmic current exactly.  A strict-right correction then gives the
continued \(\Xi\)-current.  This is a genuine pre-divisor transport and not an RH equivalence.
It establishes the singular category demanded by the no-sewing law, while making clear that a
separate zero-exclusion input is still not present.

The exact missing input can nevertheless be named without conjectural language.  A canonical
finite-aperture \(\Xi\) phase is contractive on every right-half-plane Hardy packet, and its
associated transport is four-unit nonexpansive.  Its derivative at zero aperture is the
continued anti-causal Ward column.  The derivative satisfies the four-unit scalar bound on
\(\Re w=\tfrac12\) exactly when there are no strict-right \(\Xi\) zeros.  Positive-aperture
contractivity is therefore already complete; the unresolved zero-exclusion step is
boundary-tangent control across the source-type fracture.

Independent geometric prime valuations are classical in probabilistic number theory
\cite{ArratiaBarbourTavare1997,ArratiaBarbourTavare2003}, Hellinger tensorization and Kakutani's
dichotomy are classical, and the Euler and Mertens products are classical.  Weighted-prime
test-function limits of the form used to obtain \(\Phi\) are also standard; for example,
Jana's Lemma~5.2 treats the same Mertens--Stieltjes mechanism
\cite{Jana2021}.  Negligible nets and their preservation under moderate multiplication are
classical in Colombeau theory; the same ideal mechanism and its functoriality for continuous
linear maps are reviewed in \cite{KhelifScarpalezos2021}.  We therefore claim no novelty for
the weighted-prime limit, the profile viewed in isolation, or the abstract ideal argument by
itself.  The new theorem package is the exact reflected Hellinger defect and Hellinger--Wick
reciprocity, the realization of opposite negligible/finite types by reflected zeta exponents
with a sharp essential sewing cost, and the signed pre-divisor carrier forced to cross that
fracture through a singular divisor grade.

\section{The prime-valuation product source}\label{sec:source}

Let \(\cP\) denote the rational primes.  For \(\alpha>0\), define a probability measure on
\(\N_0\) by
\[
 \mu_{p,\alpha}(k)=(1-p^{-\alpha})p^{-\alpha k}.
 \tag{2.1}\label{eq:mu-alpha}
\]
This is the geometric law of a prime valuation.  At \(\alpha=1\), the independent random
integer
\[
 M_N=\prod_{p\leq N}p^{V_p},
 \qquad V_p\sim\mu_{p,1},
 \tag{2.2}\label{eq:random-integer}
\]
has reciprocal mass on the integers supported by primes below \(N\):
\[
 \Pr(M_N=n)=\frac{a_N}{n}.
 \tag{2.3}\label{eq:reciprocal-mass}
\]
This model and its relation to prime factorization are standard
\cite{ArratiaBarbourTavare1997}.

We center the exponent at \(\alpha=1\) by writing
\[
 \alpha(c)=1-2c,\qquad |c|<\frac12,
 \tag{2.4}\label{eq:alpha-c}
\]
and abbreviate
\[
 \mu_{p,c}=\mu_{p,\alpha(c)}.
\]
Reflection \(c\mapsto-c\) sends \(\alpha(c)\) to \(2-\alpha(c)\), the radial exponent paired
about \(1\).

\begin{definition}[Finite-prime Hellinger sector]\label{def:hellinger}
For \(p\in\cP\), set
\[
 h_p(c)=\sum_{k\geq0}\sqrt{\mu_{p,0}(k)\mu_{p,c}(k)}.
 \tag{2.5}\label{eq:local-affinity-definition}
\]
For \(N\geq2\), define
\[
 A_N(c)=\prod_{p\leq N}h_p(c),\qquad H_N(c)=A_N(c)^2.
 \tag{2.6}\label{eq:global-affinity}
\]
\end{definition}

The local sum is elementary but decisive.

\begin{lemma}[Exact local affinity]\label{lem:local-affinity}
Put
\[
 q_p=p^{-1},\qquad
 r_p(c)=p^{-(1-2c)},\qquad
 u_p(c)=p^{-(1-c)}.
 \tag{2.7}\label{eq:qru}
\]
Then \(u_p(c)^2=q_pr_p(c)\), and
\[
 h_p(c)=
 \frac{\sqrt{(1-q_p)(1-r_p(c))}}{1-u_p(c)},
 \tag{2.8}\label{eq:local-affinity}
\]
\[
 h_p(c)^2=
 \frac{(1-q_p)(1-r_p(c))}{(1-u_p(c))^2}.
 \tag{2.9}\label{eq:local-squared-affinity}
\]
\end{lemma}

\begin{proof}
Substitution into \eqref{eq:local-affinity-definition} gives
\[
 \sqrt{(1-q_p)(1-r_p)}
 \sum_{k\geq0}(q_pr_p)^{k/2}.
\]
Since \((q_pr_p)^{1/2}=u_p<1\), the geometric series is
\((1-u_p)^{-1}\).  Squaring gives \eqref{eq:local-squared-affinity}.
\end{proof}

\begin{lemma}[Local order inequalities]\label{lem:local-order}
For \(0\leq q,r\leq1\), \(0\leq u<1\), and \(u^2=qr\), put
\[
 D(q,r,u)=\frac{(1-q)(1-r)}{(1-u)^2}.
\]
Then:
\begin{enumerate}[label=\textup{(\roman*)}]
\item if \(r\leq u\), then \(1-q\leq D(q,r,u)\);
\item if \(u\leq\tfrac12\), then
\[
 \frac1{(1-u)^2}\leq1+6u\leq e^{6u};
 \tag{2.10}\label{eq:exp-local-bound}
\]
\item if \(8u\leq3r\), then \(D(q,r,u)\leq1-r/4\);
\item if, in addition, \(4Lq\leq r\) for \(L\in\N\), then
\[
 D(q,r,u)\leq(1-q)^L.
 \tag{2.11}\label{eq:wold-power-local}
\]
\end{enumerate}
\end{lemma}

\begin{proof}
For (i), after multiplying by the positive denominator it is enough to show
\(1-r\geq(1-u)^2\), or \(2u-u^2\geq r\).  Since \(u^2=qr\leq r\) and
\(r\leq u\), one has \(2u-u^2\geq u\geq r\).

For (ii), multiplication by \((1-u)^2\) reduces the first inequality to
\[
 u(2u-1)(3u-4)\geq0,
\]
which holds on \(0\leq u\leq\tfrac12\); the second is \(1+x\leq e^x\).

For (iii), the difference between the right and left sides after multiplication by
\((1-u)^2\) is
\[
 q+\frac{3r-8u}{4}+\frac{ru(2-u)}4,
\]
which is nonnegative under the displayed hypotheses.  Finally,
\[
 D\leq1-\frac r4\leq1-Lq\leq(1-q)^L,
\]
where the last inequality is Bernoulli's inequality.
\end{proof}

\section{The universal Hellinger reflection law}\label{sec:phase-law}

\subsection{Critical window: Hellinger--Wick reciprocity}

For fixed \(c\in\R\), put \(\varepsilon_N=c/\log N\) and define
\[
 \mathfrak h_{p,N}(c)=
 \sum_{k\geq0}
 \sqrt{\mu_{p,\varepsilon_N}(k)\mu_{p,-\varepsilon_N}(k)},
 \qquad
 \mathcal R_N(c)=\prod_{p\leq N}\mathfrak h_{p,N}(c).
 \tag{3.1}\label{eq:reflected-critical-affinity}
\]
Also set
\[
 \Phi(c)=4\int_0^1\frac{\sinh^2(cu)}u\,du,
 \tag{3.2}\label{eq:universal-profile}
\]
where the integrand is continuously extended by \(0\) at \(u=0\), and let
\[
 K_N(z)=\sum_{p\leq N}\frac{2\{\cosh(z\log p)-1\}}p .
 \tag{3.3}\label{eq:prime-wick-cumulant}
\]
For real \(z\), \(K_N(z)\) is the cumulant of the finite symmetric Poisson prime-jump
field with atoms \(\pm\log p\), each of intensity \(1/p\).

\begin{maintheorem}[Universal critical-window law]\label{thm:critical-window}
For every \(c\in\R\):
\begin{enumerate}[label=\textup{(\roman*)}]
\item the exact local identity is
\[
 \mathfrak h_{p,N}(c)^2
 =1-x_{p,N}(c),\qquad
 x_{p,N}(c)=
 \frac{4p^{-1}\sinh^2(c\log p/\log N)}{(1-p^{-1})^2};
 \tag{3.4}\label{eq:exact-reflected-local-defect}
\]
\item the additive action and exact Hellinger defect have the same limit:
\[
 \sum_{p\leq N}\frac{4\sinh^2(c\log p/\log N)}p
 \longrightarrow\Phi(c),\qquad
 \sum_{p\leq N}x_{p,N}(c)\longrightarrow\Phi(c);
 \tag{3.5}\label{eq:additive-profile-limits}
\]
\item the multiplicative affinity satisfies
\[
 \mathcal R_N(c)^2\longrightarrow e^{-\Phi(c)};
 \tag{3.6}\label{eq:critical-hellinger-limit}
\]
\item the prime-log cumulant is asymptotically reciprocal to the Hellinger partition:
\[
 K_N(2c/\log N)\longrightarrow\Phi(c),\qquad
 \mathcal R_N(c)^2e^{K_N(2c/\log N)}\longrightarrow1.
 \tag{3.7}\label{eq:hellinger-wick-reciprocity}
\]
\end{enumerate}
\end{maintheorem}

\begin{proof}
The geometric sum in Lemma~\ref{lem:local-affinity}, now applied to the pair
\(\varepsilon_N,-\varepsilon_N\), has cross ratio \(p^{-1}\).  Hence
\[
 \mathfrak h_{p,N}(c)^2
 =
 \frac{(1-p^{-1+2\varepsilon_N})(1-p^{-1-2\varepsilon_N})}
      {(1-p^{-1})^2}.
\]
Subtracting this expression from \(1\), and using
\(e^{2y}+e^{-2y}-2=4\sinh^2y\), proves
\eqref{eq:exact-reflected-local-defect}.

For the limit, put \(L=\log N\) and introduce the finite measure
\[
 \nu_N=\frac1L\sum_{p\leq N}\frac{\log p}{p}\,
 \delta_{\log p/L}
\]
on \([0,1]\).  Mertens's first theorem in prime form
\[
 \sum_{p\leq x}\frac{\log p}{p}=\log x+O(1)
\]
implies \(\nu_N\Rightarrow du\) on \([0,1]\)
\cite{Mertens1874}.  The function
\[
 g_c(u)=
 \begin{cases}
 4\sinh^2(cu)/u,&u>0,\\
 0,&u=0,
 \end{cases}
\]
is continuous.  Therefore
\[
 \sum_{p\leq N}\frac{4\sinh^2(c\log p/L)}p
 =\int_0^1g_c(u)\,d\nu_N(u)
 \longrightarrow\int_0^1g_c(u)\,du=\Phi(c).
\]

The vacuum denominator changes no mass in this window.  For every prime \(p\),
\[
 0\leq\frac1{(1-p^{-1})^2}-1\leq\frac8p,
\]
and, for \(0\leq u\leq1\),
\[
 \sinh^2(cu)\leq c^2u^2\cosh^2|c|.
\]
Consequently
\[
 0\leq
 \sum_{p\leq N}x_{p,N}(c)-
 \sum_{p\leq N}\frac{4\sinh^2(c\log p/L)}p
 \leq
 \frac{32c^2\cosh^2|c|}{L^2}
 \sum_{p\leq N}\left(\frac{\log p}{p}\right)^2.
\]
The series on the right converges, proving the second limit in
\eqref{eq:additive-profile-limits}.

Moreover \((\log p)^2/p\leq4\) and the same estimates give
\[
 0\leq x_{p,N}(c)\leq
 \frac{80c^2\cosh^2|c|}{L^2}
\]
uniformly for \(p\leq N\).  Thus \(\max_{p\leq N}x_{p,N}(c)\to0\).  The elementary
bound
\[
 0\leq-\log(1-x)-x\leq2x^2
 \qquad(0\leq x\leq\tfrac12)
\]
then yields
\[
 \log\mathcal R_N(c)^2
 =\sum_{p\leq N}\log(1-x_{p,N}(c))
 =-\sum_{p\leq N}x_{p,N}(c)+o(1).
\]
This proves \eqref{eq:critical-hellinger-limit}.  Finally,
\[
 K_N(2c/L)
 =\sum_{p\leq N}\frac{4\sinh^2(c\log p/L)}p
\]
exactly, so \eqref{eq:hellinger-wick-reciprocity} follows from the two limits already proved.
\end{proof}

\begin{corollary}[One action, two representations]\label{cor:one-action}
The profile \(\Phi\) is even, real analytic, vanishes only at \(0\), and
\[
 \Phi(c)=2c^2+\frac13c^4+O(c^6)\qquad(c\to0).
\]
It is simultaneously the limiting L\'evy cumulant of the symmetric prime-log field and minus
the logarithm of the reflected valuation affinity.
\end{corollary}

\begin{proof}
All assertions follow from \eqref{eq:universal-profile}, the power series for \(\sinh\), and
Theorem~\ref{thm:critical-window}.
\end{proof}

\subsection{Fixed displacement: oriented Rees phases}

\begin{definition}[Rees type]\label{def:rees-type}
Let \(a_N>0\) with \(a_N\to0\).  A scalar or normed-space-valued germ \(x_N\) is
\emph{flat beyond every Rees order} if
\[
 \|x_N\|=o(a_N^m)
 \qquad\text{for every }m\in\N.
 \tag{3.8}\label{eq:all-orders-flat}
\]
The set of such germs is denoted \(\cF_\infty(a)\).  Two nonnegative germs have the
same finite Rees type if they are comparable up to positive constants.
\end{definition}

\begin{maintheorem}\label{thm:main}
Let \(A_N(c)\), \(H_N(c)\), and \(a_N\) be given by
\eqref{eq:global-affinity} and \eqref{eq:fixed-sector-intro}.  For \(|c|<\tfrac12\):
\begin{enumerate}[label=\textup{(\roman*)}]
\item If \(c<0\), then
\[
 a_N\leq H_N(c)\leq C(c)a_N
 \tag{3.9}\label{eq:negative-comparison}
\]
for all \(N\), where
\[
 C(c)=\exp\!\left(6\sum_{n\geq1}n^{-(1-c)}\right)<\infty.
 \tag{3.10}\label{eq:correction-constant}
\]
Consequently
\[
 H_N(c)=\Thetaas(a_N),\qquad
 A_N(c)=\Thetaas(a_N^{1/2}).
 \tag{3.11}\label{eq:negative-half-normal}
\]
More sharply, with \(Z_N\) from \eqref{eq:partition-curvature-intro},
\[
 H_N(c)=
 \frac{Z_N(1-c)^2}{Z_N(1)Z_N(1-2c)}
 \tag{3.12}\label{eq:exact-euler-hellinger}
\]
for every \(c\), and
\[
 \frac{H_N(c)}{a_N}\longrightarrow
 \frac{\zeta(1-c)^2}{\zeta(1-2c)}
 \qquad(c<0).
 \tag{3.13}\label{eq:exact-negative-limit}
\]
\item If \(c=0\), then \(H_N(0)=A_N(0)=1\) for every \(N\).
\item If \(c>0\), then
\[
 H(c)\in\cF_\infty(a),\qquad A(c)\in\cF_\infty(a).
 \tag{3.14}\label{eq:positive-flat}
\]
\item For every \(c\neq0\), reflection exchanges distinct types:
\[
 c\mapsto-c:
 \quad
 \Thetaas(a_N^{1/2})
 \longleftrightarrow
 \cF_\infty(a).
 \tag{3.15}\label{eq:type-fracture}
\]
\end{enumerate}
\end{maintheorem}

\begin{proof}
Suppose first that \(c<0\).  Then
\[
 r_p(c)\leq u_p(c)\leq q_p\leq\frac12.
\]
Lemma~\ref{lem:local-order}(i) gives \(h_p(c)^2\geq1-q_p\), hence
\[
 H_N(c)\geq\prod_{p\leq N}(1-q_p)=a_N.
\]
For the upper bound, \eqref{eq:local-squared-affinity} and
Lemma~\ref{lem:local-order}(ii) give
\[
 h_p(c)^2
 \leq(1-q_p)\exp(6u_p(c)).
\]
Therefore
\[
 H_N(c)
 \leq a_N\exp\!\left(6\sum_{p\leq N}p^{-(1-c)}\right)
 \leq C(c)a_N.
\]
The series in \eqref{eq:correction-constant} converges because \(1-c>1\).
Taking square roots proves \eqref{eq:negative-half-normal}.

Independently of sign, multiplying the local identity
\eqref{eq:local-squared-affinity} gives
\[
 H_N(c)=
 \prod_{p\leq N}
 \frac{(1-p^{-1})(1-p^{-(1-2c)})}
      {(1-p^{-(1-c)})^2}
 =
 \frac{Z_N(1-c)^2}{Z_N(1)Z_N(1-2c)}.
\]
Since \(a_N=Z_N(1)^{-1}\), for \(c<0\) both remaining Euler products converge
absolutely:
\[
 \frac{H_N(c)}{a_N}
 =\frac{Z_N(1-c)^2}{Z_N(1-2c)}
 \longrightarrow
 \frac{\zeta(1-c)^2}{\zeta(1-2c)}.
\]
This proves \eqref{eq:exact-euler-hellinger}--\eqref{eq:exact-negative-limit}.

If \(c=0\), then \(q_p=r_p(0)=u_p(0)\), so every local affinity is \(1\).

Now let \(c>0\) and fix \(L\in\N\).  Since
\[
 \frac{u_p(c)}{r_p(c)}=p^{-c}\longrightarrow0,
 \qquad
 \frac{q_p}{r_p(c)}=p^{-2c}\longrightarrow0,
\]
there is \(P=P(c,L)\) such that, for every prime \(p\geq P\),
\[
 8u_p(c)\leq3r_p(c),\qquad4Lq_p\leq r_p(c).
\]
Lemma~\ref{lem:local-order}(iv) yields
\[
 h_p(c)^2\leq(1-q_p)^L\qquad(p\geq P).
\]
The omitted finite prefix contributes a constant \(C_{c,L}\), so
\[
 H_N(c)\leq C_{c,L}a_N^L.
\]
Given \(m\), choose \(L=m+1\).  Since \(a_N\to0\), this implies
\(H_N(c)=o(a_N^m)\).  Applying the same argument to the square, with exponent
\(2m+1\), gives \(A_N(c)=o(a_N^m)\).  Reflection now gives
\eqref{eq:type-fracture}.
\end{proof}

The theorem reveals a genuine chirality.  Both off-critical faces approach zero, but their
orders belong to different asymptotic categories.  The negative face differs from the vacuum
half-density only by a bounded nonzero correction product.  The positive face is invisible in
every finite associated grade.

\begin{corollary}[Kakutani singularity with recovered orientation]\label{cor:kakutani}
Let
\[
 \boldsymbol\mu_c=\bigotimes_{p\in\cP}\mu_{p,c}
\]
be the product measure on \(\N_0^\cP\).  For every \(c\neq0\) with
\(|c|<\tfrac12\),
\[
 \boldsymbol\mu_c\perp\boldsymbol\mu_0.
 \tag{3.16}\label{eq:kakutani-singular}
\]
The ordinary equivalence--singularity dichotomy gives the same answer on both faces, while the
finite-product rate determines the sign of \(c\) by
\eqref{eq:type-fracture}.
\end{corollary}

\begin{proof}
Kakutani's theorem identifies the Hellinger integral of the product measures with
\(\lim_N A_N(c)\).  For \(c<0\), this limit is zero because \(a_N\to0\).  For \(c>0\),
it is zero by \eqref{eq:positive-flat}.  The rate classification is
Theorem~\ref{thm:main}.
\end{proof}

\begin{remark}[Why the infinite product space matters]\label{rem:infinite-product}
At the critical exponent, infinitely many primes are occupied almost surely, so the product
measure is not a probability law on finite integers.  It remains a perfectly well-defined
probability measure on the full valuation space \(\N_0^\cP\).  The critical boundary does not
destroy the source; it changes its support category.
\end{remark}

\section{All-orders type superselection under finite-Rees transport}\label{sec:no-sewing}

The phase distinction is stable under a much larger class than finite matrices.

\begin{definition}[Finite-Rees-moderate transport]\label{def:moderate}
Let \(E,F\) be real normed spaces.  A family of bounded linear maps
\(R_N:E\to F\) has finite Rees order if, for some \(K\in\N\),
\[
 \|R_N\|=O(a_N^{-K}).
 \tag{4.1}\label{eq:operator-grade}
\]
\end{definition}

\begin{theorem}[All-orders type-superselection law]\label{thm:flat-invariance}
If \(R_N:E\to F\) has finite Rees order and \(x\in\cF_\infty(a;E)\), then
\[
 (R_Nx_N)_N\in\cF_\infty(a;F).
 \tag{4.2}\label{eq:flat-invariance}
\]
This holds in arbitrary dimension.
\end{theorem}

\begin{proof}
Choose \(K\) as in \eqref{eq:operator-grade}.  For every \(m\),
\[
 \|R_Nx_N\|
 \leq\|R_N\|\,\|x_N\|
 =O(a_N^{-K})\,o(a_N^{m+K})
 =o(a_N^m).
\]
\end{proof}

\begin{theorem}[Universal essential sewing law]\label{thm:essential-sewing}
Let \(a_N>0\) with \(a_N\to0\), let \(E,F\) be arbitrary real normed spaces, and let
\(x_N\in E\), \(y_N\in F\).  Suppose \(x_N\neq0\) eventually,
\[
 x\in\cF_\infty(a;E),
 \qquad
 \|y_N\|=\Thetaas(a_N^d)
 \tag{4.3}\label{eq:abstract-types}
\]
for some finite \(d\in\N\).  If bounded linear maps \(R_N:E\to F\) satisfy
\[
 R_Nx_N=y_N
\]
at every cutoff, then, for every \(K\in\N\),
\[
 a_N^{-K}=o(\|R_N\|).
 \tag{4.4}\label{eq:essential-sewing}
\]
\end{theorem}

\begin{proof}
Fix \(K\).  All-orders flatness at order \(K+d\) and the finite-grade lower bound give
\[
 a_N^{-K}\|x_N\|=o(a_N^d),
 \qquad
 a_N^d=O(\|y_N\|),
\]
and hence \(a_N^{-K}\|x_N\|=o(\|y_N\|)\).
Exact sewing gives
\[
 \|y_N\|\leq\|R_N\|\,\|x_N\|.
\]
After division by the eventually nonzero norms, the ratio
\(a_N^{-K}/\|R_N\|\) is bounded above by
\(a_N^{-K}\|x_N\|/\|y_N\|\), which tends to zero.  This proves
\eqref{eq:essential-sewing}.
\end{proof}

\begin{corollary}[Essential prime-reflection sewing cost]\label{cor:no-sewing}
Let \(0<c<\tfrac12\), and set
\[
 S_N(c)=\frac{H_N(-c)}{H_N(c)}.
\]
Then
\[
 a_N^{-K}=o\!\left(S_N(c)\right)
 \qquad\text{for every }K\in\N.
 \tag{4.5}\label{eq:prime-essential-sewing}
\]
Every exact bounded linear sewing between normed-space realizations of the positive and negative
prime sectors has the same essential operator-norm singularity.  In particular, no
finite-Rees-moderate operator family can satisfy
\[
 R_NH_N(c)=H_N(-c)
 \qquad\text{for every }N.
 \tag{4.6}\label{eq:scalar-no-sewing}
\]
\end{corollary}

\begin{proof}
Every finite-prime Hellinger factor is strictly positive.
Theorem~\ref{thm:main} gives
\[
 H_N(c)\in\cF_\infty(a),\qquad H_N(-c)=\Thetaas(a_N).
\]
Apply Theorem~\ref{thm:essential-sewing} with \(d=1\).  The scalar quotient is \(S_N(c)\);
the operator statement and the no-sewing conclusion follow from the same theorem.
\end{proof}

\begin{corollary}[Finite packets and invertible changes of frame]\label{cor:finite-packets}
All-orders flatness of a finite vector packet is equivalent to coordinatewise all-orders
flatness.  Every finite matrix whose entries have one uniform finite Rees order preserves the
flat radical.  If the matrix and its inverse are both finite-Rees-moderate, they preserve and
reflect the radical exactly.
\end{corollary}

\begin{proof}
All norms on a finite product are equivalent, so a finite packet is flat exactly when all of its
coordinates are flat.  Each output coordinate is a finite sum of products of a finite-grade
coefficient and a flat input.  Applying the same argument to the inverse proves the last claim.
\end{proof}

The conclusion is narrower and stronger than a slogan about unbounded operators.  The reflected
faces do not have merely unequal finite grades: one is half-normal and the other belongs to the
intersection of every positive Rees grade.  That intersection is an invariant radical.
Unboundedness at one cutoff is irrelevant; the theorem distinguishes every fixed polynomial
order from growth beyond all finite Rees orders.  Likewise, infinite rank is irrelevant unless
it changes the full operator-norm type.  Theorem~\ref{thm:essential-sewing} quantifies the only
possible change: exact sewing forces the full norm to dominate every inverse Rees power.

\section{Reflection versus modular retuning}\label{sec:modular-retuning}

The no-sewing theorem would be uninformative if it merely ruled out a small class and said
nothing about how reflection escapes.  A rank-one Weyl block gives the exact positive answer.
Consider
\[
 \lambda_p(z)=\frac{pz^2-1}{p-z^2},
 \qquad p>1,\quad p-z^2\neq0.
 \tag{5.1}\label{eq:weyl-eigenvalue}
\]
Calegari and Venkatesh give an explicit two-channel level-changing scattering block
\(N(z,s)\) induced by the standard rank-one intertwiner~\cite[\S6.6.1]{CalegariVenkatesh}.
At \(s=0\), its eigenvalues are \(1\) and \(\lambda_p(z)\).  We need only the following
algebraic law of the second eigenvalue.

\begin{theorem}[Exact modular-retuning law]\label{thm:weyl-retuning}
For every \(p>1\) and \(z\in\C\) with \(p-z^2\neq0\),
\[
 |pz^2-1|^2-|p-z^2|^2
 =(p^2-1)(|z|^4-1).
 \tag{5.2}\label{eq:weyl-modulus-defect}
\]
Consequently
\[
 |\lambda_p(z)|=1\quad\Longleftrightarrow\quad |z|=1.
 \tag{5.3}\label{eq:weyl-unitary-locus}
\]
If \(z=p^c\) with \(0<c<\tfrac12\), then
\[
 p^{2c}<
 \frac{p\,p^{2c}-1}{p-p^{2c}}
 =\lambda_p(p^c).
 \tag{5.4}\label{eq:weyl-radial-cost}
\]
Finally, for every nonzero scalar gauge \(g\),
\[
 \frac{g\lambda_p(z)}{g}=\lambda_p(z);
 \tag{5.5}\label{eq:weyl-gauge-ratio}
\]
scalar normalization can move the growth between the two channels but cannot remove their
ratio.
\end{theorem}

\begin{proof}
Writing \(w=z^2\), direct expansion gives
\[
 |pw-1|^2=p^2|w|^2-2p\Re w+1,
 \qquad
 |p-w|^2=p^2-2p\Re w+|w|^2,
\]
whose difference is \eqref{eq:weyl-modulus-defect}.  Since \(p^2-1>0\), equality of the two
moduli is equivalent to \(|z|=1\), proving \eqref{eq:weyl-unitary-locus}.

For \(r=p^{2c}\), one has \(1<r<p\).  Hence
\[
 \frac{pr-1}{p-r}>r
 \quad\Longleftrightarrow\quad
 pr-1>pr-r^2
 \quad\Longleftrightarrow\quad r^2>1.
\]
This proves \eqref{eq:weyl-radial-cost}; \eqref{eq:weyl-gauge-ratio} is cancellation of a
nonzero scalar.
\end{proof}

The theorem separates two notions that scalar functional equations conflate:
\[
 \boxed{\quad
 \begin{array}{c}
 \text{source-internal reflection}\\[-1mm]
 \text{with fixed modular class}
 \end{array}
 \quad\text{versus}\quad
 \begin{array}{c}
 \text{external Weyl joining}\\[-1mm]
 \text{with modular retuning}.
 \end{array}
 \quad}
 \tag{5.6}\label{eq:reflection-categories}
\]
The first category is constrained by Theorem~\ref{thm:flat-invariance}; the second can cross the
fracture by carrying the missing growth in its complementary channel.  Criticality is the unique
locus on which the retuning cost vanishes locally.  Thus the finite-Rees result is not a no-go
for adelic reflection.  It is a conservation theorem for reflection that remains internal to
the fixed ordinary-prime source.

\section{Diagonal activation and the critical source boundary}\label{sec:activation}

The Wold scale in Theorem~\ref{thm:main} is not an external normalization.  It arises from an
exact probability disintegration of the diagonal Euler activation.

For \(\Re s>1\),
\[
 \cZ_\infty(s,s)=\prod_p(1+p^{-s})
 =\frac{\zeta(s)}{\zeta(2s)}.
 \tag{6.1}\label{eq:diagonal-activation}
\]
At a local node \(q=p^{-\sigma}\), normalize the positive coefficients by \(1+q\).
The resulting valuation law is
\[
 \nu_q(0)=\frac1{1+q},
 \qquad
 \nu_q(k)=\frac{(1-q)q^k}{1+q}\quad(k\geq1).
 \tag{6.2}\label{eq:activation-law}
\]

\begin{proposition}[Occupancy--fibre disintegration]\label{prop:disintegration}
For the law \(\nu_q\):
\begin{enumerate}[label=\textup{(\roman*)}]
\item the radical occupancy \(B=\mathbf1_{\{V>0\}}\) is Bernoulli with
\[
 \Pr(B=1)=\frac q{1+q};
 \tag{6.3}\label{eq:occupancy}
\]
\item conditional on \(B=1\), the positive valuation is geometric:
\[
 \Pr(V=k\mid B=1)=(1-q)q^{k-1},\qquad k\geq1;
 \tag{6.4}\label{eq:conditional-geometric}
\]
\item
\[
 \Pr(V\geq m)=\frac{q^m}{1+q},\qquad m\geq1;
 \tag{6.5}\label{eq:valuation-tail}
\]
\item the occupancy variance is
\[
 \operatorname{Var}(B)=\frac q{(1+q)^2}.
 \tag{6.6}\label{eq:fisher-variance}
\]
\end{enumerate}
\end{proposition}

\begin{proof}
Summing the positive geometric tail gives
\[
 \sum_{k\geq1}\frac{(1-q)q^k}{1+q}=\frac q{1+q}.
\]
Division by this mass gives \eqref{eq:conditional-geometric}; the same summation starting at
\(m\) gives \eqref{eq:valuation-tail}.  Formula \eqref{eq:fisher-variance} is the Bernoulli
variance.
\end{proof}

At the critical node \(q=1/p\), the decomposition becomes arithmetically sharp.

\begin{theorem}[Critical occupancy law]\label{thm:critical-occupancy}
At \(q=1/p\):
\begin{enumerate}[label=\textup{(\roman*)}]
\item occupancy mass is \((p+1)^{-1}\) and is not summable over primes;
\item for every fixed \(m\geq2\),
\[
 \Pr(V\geq m)=\frac1{p^{m-1}(p+1)}
\tag{6.7}\label{eq:critical-tail}
\]
is summable over primes;
\item the Fisher-whitened Gamma fibre has mass
\[
 \frac{p-1}{p(p+1)}
\tag{6.8}\label{eq:gamma-fibre}
\]
and is again not summable over primes;
\item no nonzero polynomial vanishes on infinitely many critical nodes \(1/p\).
\end{enumerate}
\end{theorem}

\begin{proof}
The first three formulas follow from
\eqref{eq:occupancy}--\eqref{eq:fisher-variance}.  The first and third terms are bounded below
by positive constant multiples of \(p^{-1}\), while
\eqref{eq:critical-tail} is bounded by \(p^{-m}\).  Euler's divergence of
\(\sum_pp^{-1}\) and convergence of the \(p\)-series for \(m\geq2\) give the summability
claims.  The nodes \(1/p\) are infinitely many distinct roots, so a polynomial vanishing on
infinitely many of them is zero.
\end{proof}

This theorem isolates the critical source:
\[
 \boxed{\begin{gathered}
 \text{prime-harmonic occupancy}
 \ \times\
 \text{summable conditional depth fibre}\\
 \xrightarrow{\text{Fisher whitening}}
 \text{prime-harmonic Gamma fibre}
 \end{gathered}}
 \tag{6.9}\label{eq:source-box}
\]
Whitening must amplify a rare event.  Its divergent recovery gauge is therefore forced by the
source law; it is not evidence that an additional bounded coordinate has been omitted.

\section{Depth--support superselection}\label{sec:superselection}

The valuation source has a second filtration independent of Rees type.  It explains why finite
differential closure cannot replace the missing singular transport.

Let \(I\) be finite, let \(A_i\) be finite pointed sets, and set
\[
 A=\prod_{i\in I}A_i.
\]
Let \(y:\cP\to A\) visit every state on every cofinite prime tail.  The anchored tensor basis
\((b_z)_{z\in A}\) of \(\C^A\) separates the exact coordinate support
\[
 \supp(z)=\{i:z_i\neq1\}.
\]
For valuation depth \(r\geq0\), form the prime-supported Dirichlet carrier
\[
 F_{r,z}(s)=\sum_{p\in\cP}b_z(y(p))p^{-(r+1)s}
\tag{7.1}\label{eq:carrier}
\]
in its half-plane of convergence.

There is an exact same-prime fusion on coefficient columns:
\[
 V_e(f)\star V_r(g)=V_{e+r+1}(fg).
\tag{7.2}\label{eq:fusion}
\]
Thus depth adds, while coordinate support unions.

\begin{theorem}[Depth--support superselection]\label{thm:superselection}
For every \(d>0\) and \(S\subseteq I\), the family
\[
 \{F_{r,z}:0\leq r<d,\ \supp(z)=S\}
\tag{7.3}\label{eq:support-sector}
\]
is differentially algebraically independent over the differential field generated by
\(\C(s)\) and its complete complement in the depth-\(d\) rectangle.  Its exact relative rank is
\[
 d\prod_{i\in S}(|A_i|-1).
\tag{7.4}\label{eq:support-rank}
\]
The complete layer \(d\) is relatively hypertranscendental over all earlier layers.  Nevertheless
fusion of two carriers in the earlier prefix can land exactly in that excluded next layer.
\end{theorem}

\begin{proof}
Tail completeness makes the coefficient patterns of the anchored basis linearly independent
after deletion of any finite prime set.  Applying successive logarithmic derivatives multiplies
the depth-\(r\) prime coefficient by powers of \((r+1)\log p\).  At any finite differential
order these form Vandermonde blocks of full rank.  Uniqueness of Dirichlet series then gives joint
differential independence of every finite depth--state rectangle over \(\C(s)\).

The exact-support sector contains
\(\prod_{i\in S}(|A_i|-1)\) anchored states, and the depth interval contributes \(d\), proving
\eqref{eq:support-rank}.  Relative independence from the complete complement follows by splitting
a jointly independent family.

Finally, \eqref{eq:fusion} sends depths \(e,r\) to \(e+r+1\).  On disjoint anchored supports,
pointwise basis multiplication reaches the union support exactly.  Taking one input in the top
prefix layer and the neutral depth-zero source produces a carrier in the next layer, which the
preceding independence statement excludes from the differential field of the complete prefix.
\end{proof}

The two filtrations are independent:
\[
 \begin{array}{c|c|c}
 \text{filtration}&\text{finite operation}&\text{invariant obstruction}\\ \hline
 \text{valuation depth}&\text{differential-rational closure}&
   \text{new HT layer}\\
 \text{interaction support}&\text{polynomial multiplication}&
   \text{support cannot exceed a union}\\
 \text{Rees type}&\text{finite-order bounded transport}&
   \text{flat radical is preserved}.
 \end{array}
\tag{7.5}\label{eq:three-filtrations}
\]
This is the source-level reason that adding finitely many derivatives, residue channels, jets,
or Hilbert coordinates cannot manufacture the reflected half-normal face.

\section{Divisor differentiation forces singular completion}\label{sec:divisor}

The no-sewing theorem identifies what regular transport cannot do.  The next result identifies
the exact singularity required at a divisor.

\begin{proposition}[Universal divisor grade drop]\label{prop:divisor-drop}
Let \(d\geq1\), \(B(t)=t^d\), and \(S(t)=dt^{d-1}\).  For \(t\neq0\),
\[
 \frac{S(t)^2}{B(t)^2}=\frac{d^2}{t^2}.
\tag{8.1}\label{eq:inverse-square}
\]
If
\[
 G(t)=
 \begin{pmatrix}
 B(t)^2&S(t)\\
 S(t)&C(t)
 \end{pmatrix}
\tag{8.2}\label{eq:divisor-gram}
\]
is positive semidefinite, then
\[
 C(t)\geq\frac{d^2}{t^2}.
\tag{8.3}\label{eq:diagonal-blowup}
\]
Consequently no uniformly bounded positive completion exists near a nonzero divisor.
\end{proposition}

\begin{proof}
Equation \eqref{eq:inverse-square} is immediate.  Positivity of a real symmetric
\(2\times2\) matrix implies nonnegative determinant:
\[
 B(t)^2C(t)-S(t)^2\geq0.
\]
Division by \(B(t)^2>0\) gives \eqref{eq:diagonal-blowup}.
\end{proof}

The conclusion is independent of the multiplicity \(d\): multiplicity changes the coefficient,
but differentiation always drops exactly one divisor grade and positivity always squares that
loss.  A bounded completely positive extension is therefore in the wrong category at every
zero.  The natural object is a divisor-graded graph or Rees correspondence in which division by
the overlap is part of the domain.

\section{An exact singular source graph to the \(\Xi\)-current}\label{sec:singular-graph}

The exceptional category forced by Proposition~\ref{prop:divisor-drop} is not empty.  The
squarefree activation supplies an exact singular graph before the divisor is selected.

Write
\[
 Z_\square(s)=\frac{\zeta(s)}{\zeta(2s)},\qquad
 C_\square(s)=-\frac{d}{ds}\log Z_\square(s).
\tag{9.1}\label{eq:squarefree-current}
\]
For \(\Re\sigma>1\), \(\varepsilon\neq0\), and
\(s=\sigma+u\varepsilon\) with \(\Re s>1\), define the moving overlap
\[
 B_{\sigma,\varepsilon}(u)=
 \frac{Z_\square(\sigma+u\varepsilon)}{Z_\square(\sigma)}
\tag{9.2}\label{eq:moving-overlap}
\]
and the unnormalized pointed Fisher numerator
\[
 S_{\sigma,\varepsilon}(u)=
 2\left(
 \frac{d}{du}B_{\sigma,\varepsilon}(u)
 +
 \varepsilon C_\square(\sigma)B_{\sigma,\varepsilon}(u)
 \right).
\tag{9.3}\label{eq:fisher-numerator}
\]

\begin{proposition}[Exact overlap graph]\label{prop:overlap-graph}
On the Euler domain,
\[
 S_{\sigma,\varepsilon}(u)
 =2\varepsilon B_{\sigma,\varepsilon}(u)
 \bigl(C_\square(\sigma)-C_\square(s)\bigr).
\tag{9.4}\label{eq:score-identity}
\]
Hence the singular graph transform
\[
 \mathfrak G_{\sigma,\varepsilon}[S](u)=
 2C_\square(\sigma)
 -\frac{S(u)}{\varepsilon B_{\sigma,\varepsilon}(u)}
 -2\frac{d}{ds}\log A_\infty(s),
\tag{9.5}\label{eq:graph-transform}
\]
where
\[
 A_\infty(s)=\frac12s(s-1)\pi^{-s/2}\Gamma(s/2),
\]
satisfies
\[
 \mathfrak G_{\sigma,\varepsilon}
 [S_{\sigma,\varepsilon}](u)
 =
 2\left(C_\square(s)-\frac{d}{ds}\log A_\infty(s)\right).
\tag{9.6}\label{eq:completed-squarefree}
\]
Adding the strict-right correction \(4(-\zeta'/\zeta)(2s)\) yields the completed
\(\Xi\)-logarithmic current exactly.
\end{proposition}

\begin{proof}
Logarithmic differentiation of \eqref{eq:moving-overlap} gives
\[
 B'_{\sigma,\varepsilon}(u)
 =-\varepsilon B_{\sigma,\varepsilon}(u)C_\square(s).
\]
Substitution into \eqref{eq:fisher-numerator} proves \eqref{eq:score-identity}; substitution of
\eqref{eq:score-identity} into \eqref{eq:graph-transform} proves
\eqref{eq:completed-squarefree}.

Since
\[
 C_\square(s)
 =-\frac{\zeta'(s)}{\zeta(s)}
 +2\frac{\zeta'(2s)}{\zeta(2s)},
\]
the correction \(4(-\zeta'/\zeta)(2s)\) cancels the doubled squarefree subtraction and leaves
the completed \(\Xi\)-current with the stated normalization.
\end{proof}

\begin{corollary}[The singularity is source-forced]\label{cor:source-forced}
If \(Z_\square(s)\) has a zero of order \(d\), then the overlap
\(B_{\sigma,\varepsilon}\) has order \(d\) while the numerator
\(S_{\sigma,\varepsilon}\) has order \(d-1\).  The quotient in
\eqref{eq:graph-transform} has a first-order pole, and every positive Gram realization of it
requires the inverse-square diagonal growth of Proposition~\ref{prop:divisor-drop}.
\end{corollary}

\begin{proof}
The derivative of a function with a zero of order \(d\) has order \(d-1\).
Equation \eqref{eq:fisher-numerator} adds a term divisible by the overlap, so it does not cancel
the leading derivative grade.  The remaining statements are
Proposition~\ref{prop:divisor-drop}.
\end{proof}

The graph is the first positive survivor of the bounded no-go.  It divides before scalar
continuation loses the source grade, incorporates the Gamma completion, and reaches the
\(\Xi\)-current without enumerating zeros.  It is not a bounded dagger map and is not claimed to
be one.

\section{Four-unit contractivity and the boundary-tangent fracture}
\label{sec:ward-boundary}

The critical profile \(\Phi\) is even.  It measures reflection cost but, like every squared
distance, cannot retain the sign of the tangent.  The relevant HT-geometric scalarization is
instead a polarization.

Let \(\mathcal H\) be a complex Hilbert space, let \(e,T\) be bounded operators, and use the
commutator convention \([T,e]=Te-eT\).  For \(\psi,\chi\in\mathcal H\), define the pointed
Fisher connection
\[
 F_{\psi,\chi}=\frac1{\sqrt2}
 \bigl(|\chi\rangle\langle\psi|+|\psi\rangle\langle\chi|\bigr)
 \tag{10.1}\label{eq:pointed-fisher}
\]
and the cutoff Ward pairing
\[
 \mathfrak W_{e,T}(\psi,\chi)=
 \operatorname{Tr}\bigl([F_{\psi,\chi},e]^*[T,e]\bigr).
 \tag{10.2}\label{eq:pointed-ward}
\]
The trace is well defined although \(e\) and \(T\) need not be Hilbert--Schmidt, because
\(F_{\psi,\chi}\) has rank at most two.

\begin{theorem}[Signed Fisher--Ward transgression]\label{thm:signed-transgression}
For arbitrary bounded \(e,T\):
\[
 \mathfrak W_{e,T}(\psi,-\chi)
 =-\mathfrak W_{e,T}(\psi,\chi),
 \tag{10.3}\label{eq:ward-odd}
\]
and
\[
 \begin{aligned}
 \mathfrak W_{e,T}(\psi,\chi)
 &=\operatorname{Tr}\!\left(
 F_{\psi,\chi}\,[[T,e],e^*]\right)\\
 &=\frac1{\sqrt2}\left(
 \langle\psi,[[T,e],e^*]\chi\rangle+
 \langle\chi,[[T,e],e^*]\psi\rangle\right).
 \end{aligned}
 \tag{10.4}\label{eq:signed-transgression}
\]
If \(e=e^*\), the curvature is the ordinary double commutator \([[T,e],e]\).
\end{theorem}

\begin{proof}
The connection is linear in \(\chi\), hence changes sign under
\(\chi\mapsto-\chi\); so does \eqref{eq:pointed-ward}.  Expanding both commutators,
passing bounded factors cyclically across the finite-rank connection, and regrouping gives
\[
 \operatorname{Tr}\bigl([F,e]^*[T,e]\bigr)
 =\operatorname{Tr}\bigl(F[[T,e],e^*]\bigr).
\]
The trace of \(|x\rangle\langle y|B\) is \(\langle y,Bx\rangle\).  Applying this
rank-one identity to the two summands of \(F_{\psi,\chi}\) proves the second line.
\end{proof}

Theorem~\ref{thm:signed-transgression} isolates what squaring had removed.  The universal
Hellinger--Wick action supplies the size of the prime tangent; its pointed Fisher polarization
supplies the sign.  Both are defined before divisor selection.  A source-internal RH route must
therefore transport the pointed connection, or an equivalent odd object, through the singular
graded graph.  Transporting only \(\Phi\), a norm, an affinity, or a divisor count cannot recover
orientation after the fact.

The singular current has an exact Hardy-space consumer.  Let \(\cD_R^+\) be the upper
Schwarz representatives of strict-right centered \(\Xi\) zeros.  For \(\Re w>0\) and
\(\alpha\geq0\), define the canonical-product phase
\[
 \Phi_\alpha(w)=
 \prod_{\lambda\in\cD_R^+}
 \frac{w+\lambda-\alpha}{w+\lambda+\alpha}\,
 \frac{w+\bar\lambda-\alpha}{w+\bar\lambda+\alpha},
 \qquad
 V_\alpha(w)=1-\Phi_\alpha(w),
 \tag{10.5}\label{eq:ward-phase}
\]
where the product is taken in the canonical Schwarz-paired normalization.  Each factor is a
right-half-plane Cayley contraction.

\begin{theorem}[Finite aperture versus boundary tangent]\label{thm:ward-fracture}
Let \(K(w,z)=(w+\bar z)^{-1}\) be the right-half-plane Hardy kernel.
\begin{enumerate}[label=\textup{(\roman*)}]
\item For every \(\alpha\geq0\), \(|\Phi_\alpha(w)|\leq1\) and
\[
 |V_\alpha(w)|\leq2.
 \tag{10.6}\label{eq:ward-pointwise}
\]
\item On every finite packet \(w_1,\dots,w_n\) in the right half-plane,
\[
 \left[
 4K(w_i,w_j)
 -V_\alpha(w_i)K(w_i,w_j)\overline{V_\alpha(w_j)}
 \right]_{i,j=1}^n
 \succeq0.
 \tag{10.7}\label{eq:ward-four-unit-pick}
\]
\item The zero-aperture derivative is the continued anti-causal Ward response:
\[
 \left.\frac{\partial}{\partial\alpha}V_\alpha(w)\right|_{\alpha=0}
 =\cR_\Xi(w).
 \tag{10.8}\label{eq:ward-derivative}
\]
\item On the fixed Euler edge,
\[
 \boxed{\quad
 \text{\rm RH}
 \quad\Longleftrightarrow\quad
 |\cR_\Xi(\tfrac12+it)|^2\leq4
 \quad\text{for every }t\in\R.
 \quad}
 \tag{10.9}\label{eq:ward-rh-consumer}
\]
If \(\lambda=a+ib\) is a strict-right centered zero, then the adapted node
\(w=\tfrac12-ib\) satisfies
\[
 |\cR_\Xi(w)|^2>4.
 \tag{10.10}\label{eq:ward-adapted-violation}
\]
\end{enumerate}
\end{theorem}

\begin{proof}
For \(\Re w,\Re\lambda,\alpha>0\),
\[
 \left|\frac{w+\lambda-\alpha}{w+\lambda+\alpha}\right|\leq1.
\]
Schwarz pairing, convergence of the canonical product, and passage to the limit give
\(|\Phi_\alpha|\leq1\).  The affine Schur transform
\((1-\Phi_\alpha)/2\) is again a Schur multiplier; its Pick matrix is positive, which is
\eqref{eq:ward-four-unit-pick}.  The pointwise bound is its one-node consequence.
Differentiating the normalized canonical-product ratio at \(\alpha=0\) gives the absolutely
convergent Schwarz-paired Cauchy response, proving \eqref{eq:ward-derivative}.

If there are no strict-right zeros, the product is empty and
\(\cR_\Xi=0\).  Conversely, let \(\lambda=a+ib\) with \(0<a<\tfrac12\).  At
\(w=\tfrac12-ib\), the selected Schwarz pair contributes real part
strictly larger than \(2\); every other pair contributes positive real part.  Hence
\(\Re\cR_\Xi(w)>2\), which implies \eqref{eq:ward-adapted-violation}.  This proves
\eqref{eq:ward-rh-consumer}.
\end{proof}

The theorem prevents a common category error.  The four-unit inequality is already true at every
positive aperture; differentiating it at the boundary does not preserve the same bound.  Under
false RH the failure is not hidden in a large packet: one zero-adapted node witnesses it.
Accordingly, additional finite-aperture positivity cannot solve the problem.  The exact missing
operation is a source-controlled passage from the finite-aperture contraction to its
zero-aperture tangent.

\section{Consequences for Riemann-hypothesis proof architecture}\label{sec:rh}

The results answer structural questions that are often conflated.

\subsection{What is the universal object behind the separate obstructions?}

It is not a scalar zeta value or a zero count.  It is the reflected valuation tangent equipped
with four pieces of structure:
\[
 \boxed{
 \begin{gathered}
 \text{even action }\Phi
 \;+\;
 \text{odd pointed Fisher polarization}\\
 {}+\;
 \text{depth--support filtration}
 \;+\;
 \text{Rees asymptotic type}
 \end{gathered}}
\tag{11.1}\label{eq:universal-source-object}
\]
The critical-window theorem identifies the first term simultaneously as Hellinger loss and
prime-log fluctuation.  The signed transgression theorem proves that the second retains tangent
orientation and couples canonically to Ward curvature.  Hypertranscendence proves that the
depth--support filtration cannot be synthesized by finite differential closure.
Type superselection proves independently that the all-orders-flat radical cannot be erased by
finite-Rees-moderate transport; exact sewing has norm beyond every inverse Rees power.  The
M\"obius--Gamma graph supplies the singular divisor-grade change, and the four-unit Ward response
is the final consumer.

This explains why the previous scalar shadows were individually insufficient.  Hellinger
affinity and prime-clock energy see magnitude but are even.  Rees type sees the outer side but
does not by itself control a boundary derivative.  Differential independence prevents a finite
jet surrogate.  The pointed Fisher--Ward pairing is the first object in the chain that is both
source-defined and odd.  The remaining RH problem is now a specific compatibility estimate:
carry that odd, graded source tangent through the exact singular graph while retaining the
four-unit boundary norm.

\subsection{Where is the side of the critical line stored?}

It is stored in the rate class of the prime-valuation Hellinger sector.  The sign of \(c\) is
equivalent to a source-side asymptotic distinction:
\[
 c<0\iff A_N(c)=\Thetaas(a_N^{1/2}),
 \qquad
 c>0\iff A(c)\in\cF_\infty(a).
\tag{11.2}\label{eq:orientation-answer}
\]
This information exists before a zero is chosen.

\subsection{Why does ordinary product-measure singularity not orient reflection?}

Kakutani's theorem records only whether \(\lim_NA_N(c)\) is zero.  It is zero for both
signs.  The dichotomy is therefore a quotient of the finer type space:
\[
 \{\text{half-normal},\text{all-orders-flat}\}
 \longrightarrow
 \{\text{singular}\}.
\tag{11.3}\label{eq:kakutani-quotient}
\]
The missing sign is lost by this quotient, not absent from the Euler source.

\subsection{Why do bounded Hilbert and finite-jet completions fail to sew the pair?}

They preserve \(\cF_\infty(a)\) whenever their full norm has fixed finite Rees order.
This statement is independent of dimension and of the chosen coordinate frame.  It applies to a
large class of bounded Hilbert-space, finite-matrix, and polynomially renormalized realizations.
It does \emph{not} rule out every conceivable Hilbert--P\'olya or trace-formula construction; it
rules out exactly those whose source reflection is finite-Rees-moderate.

\subsection{How can reflection cross the fracture?}

By paying modular drift in a complementary channel.  Theorem~\ref{thm:weyl-retuning} gives an
explicit rank-one model: the reflection block is unitary precisely on \(|z|=1\), while the radial
parameter \(z=p^c\) forces a local cost larger than \(p^{2c}\) for \(c>0\).  Scalar gauge cannot
remove that channel ratio.  This answers the escape question positively and prevents the
finite-Rees obstruction from being misread as a no-go for the standard adelic Weyl mechanism.

\subsection{What object survives the obstruction?}

The overlap graph \eqref{eq:graph-transform} survives.  Its division by
\(B_{\sigma,\varepsilon}\) is precisely the divisor-grade singularity demanded by
Proposition~\ref{prop:divisor-drop}.  The activation law shows why Fisher/Gamma whitening is
unbounded, and Theorem~\ref{thm:superselection} shows why no finite differential field of
valuation layers can replace the graph.

\subsection{Where does zero exclusion remain?}

At the zero-aperture derivative.  Theorem~\ref{thm:ward-fracture} proves four-unit
nonexpansiveness at every finite aperture and identifies the derivative with the continued Ward
response.  The sharp derivative bound is RH-complete, and one off-line pair violates it at one
adapted Euler-edge node.  Thus the remaining task is neither another positivity construction nor
another scalar equivalence: it is a source-derived tangent estimate that survives the
divisor-grade and reflection-type fractures.

These answers combine into one conservation law:

\begin{theorem}[Source-type conservation law]\label{thm:conservation}
Every transport whose source action is represented by a finite valuation/residue packet followed
by a bounded operator family of fixed finite Rees order preserves the all-orders-flat radical.
It cannot convert a positive reflected prime-valuation sector into its negative reflected
sector.  Independently, no finite differential prefix of the valuation grid is closed under
the exact source fusion.  A reflection architecture must therefore declare which of two
mathematically distinct operations it performs:
\begin{enumerate}[label=\textup{(\roman*)}]
\item \emph{source-internal transport}, which retains the ordinary-prime modular class and must
use a divisor-graded correspondence whose exact reflected sewing cost is beyond every finite
Rees order; or
\item \emph{external modular retuning}, which can cross the fracture by carrying the drift in a
complementary Weyl channel.
\end{enumerate}
The squarefree M\"obius--Gamma graph realizes the singular internal category on the Euler domain
and reconstructs the continued \(\Xi\)-current.  The rank-one Weyl block realizes the external
retuning category.  The finite-aperture Ward transport is contractive, but its RH-complete
boundary tangent is not controlled by either fact alone.
\end{theorem}

\begin{proof}
Finite channel enlargement preserves flatness coordinatewise by
Corollary~\ref{cor:finite-packets}, and
Theorem~\ref{thm:flat-invariance} handles the bounded operator stage.  The reflected sectors have
opposite types by Theorem~\ref{thm:main}, so such a regular transport cannot sew them.
Theorem~\ref{thm:superselection} separately excludes fusion closure inside every finite
differential prefix.  Proposition~\ref{prop:overlap-graph} supplies the stated singular
realization, and Theorem~\ref{thm:weyl-retuning} supplies the complementary modular-retuning
realization.  Theorem~\ref{thm:ward-fracture} identifies the exact boundary consumer.
\end{proof}

\section{Scope and logical boundary}\label{sec:scope}

The paper proves an unconditional source law, not the Riemann hypothesis.  The distinction is
exact:

\begin{enumerate}[label=\textup{(\arabic*)}]
\item The reflected Hellinger critical-window law is a theorem for the actual rational primes:
its exact profile is \(\Phi(c)\), and its product is asymptotically reciprocal to the
prime-log Poisson/Wick cumulant.
\item The Hellinger--Rees outer trichotomy is exact for the reflected zeta exponents; its
negative-face normalization is the zeta ratio
\(\zeta(1-c)^2/\zeta(1-2c)\).
\item The all-orders type-superselection and essential-sewing theorems are unconditional,
arithmetic-independent, and arbitrary-dimensional.  The prime-valuation sectors realize their
flat and finite-grade inputs exactly.
\item The pointed Fisher--Ward pairing is an unconditional finite-rank trace, is odd in the
source tangent, and transgresses exactly to double-cutoff curvature.
\item The Weyl modulus identity and its unitary locus are exact.  Their role is to exhibit
modular retuning as a positive escape, not to infer RH from local unitarity.
\item The occupancy/fibre disintegration, differential superselection, divisor grade drop, and
singular squarefree-to-\(\Xi\) graph are exact identities or theorems.
\item Finite-aperture four-unit nonexpansiveness, identification of its boundary derivative with
the Ward response, and the Euler-edge four-unit zero-exclusion criterion are exact.
\item No statement here proves that every nontrivial zero must be represented by a
finite-Rees-moderate sewing operator.  Indeed the singular graph shows why such an assumption
would be unnatural.
\item The graph reconstructs the completed logarithmic current on its safe source domain.  It
does not prove that the continued current has no off-line poles.
\end{enumerate}

Thus the contribution is neither a decorative RH equivalence nor a hidden conditional proof.
It changes and sharply localizes the proof target.  The source supplies a universal even action,
an odd orientation carrier, a depth--support filtration, and an asymptotic type.  Standard
external reflection can cross the type fracture only by modular retuning; internal transport
must retain the pointed tangent through a singular divisor grade.  The remaining zero-exclusion
statement is precisely the
four-unit estimate for the zero-aperture Ward tangent.  Theorems
\ref{thm:critical-window}, \ref{thm:signed-transgression}, \ref{thm:main},
\ref{thm:flat-invariance}, and \ref{thm:weyl-retuning} classify the source data and reflection
categories, Proposition~\ref{prop:overlap-graph} gives an explicit internal singular graph, and
Theorem~\ref{thm:ward-fracture} gives the exact analytic consumer.

\section{Conclusion}\label{sec:conclusion}

The universal theorem is an asymptotic superselection law: an all-orders-flat germ and a
nonzero finite-grade germ cannot be connected by finite-order bounded transport, in any
dimension, and every exact connection has an essential scale singularity.  Prime-valuation
reflection supplies a canonical arithmetic realization of its two asymptotic species.

At the shrinking critical window, reflection is governed by the finite-size action
\[
 \Phi(c)=4\int_0^1\frac{\sinh^2(cu)}u\,du.
\]
Its exponential is the limiting reflected prime-valuation affinity, while \(\Phi\) itself is
the limiting cumulant of the symmetric prime-log field.  The reciprocity
\(\mathcal R_N(c)^2e^{K_N(2c/\log N)}\to1\) identifies information loss and prime-clock
fluctuation as two representations of one source law.

Because this action is even, it cannot orient the critical tangent.  The complementary universal
law is the pointed Fisher--Ward transgression: its trace is odd in the source tangent and equals
the pairing of that tangent with double-cutoff curvature.  This is the pre-divisor orientation
carrier that scalar affinities, squared normal coordinates, and divisor counts erase.

At fixed displacement the same valuation family changes asymptotic type.  The negative face has
the canonical half-normal Wold grade, the critical face has unit affinity, and the positive face
is flat beyond every finite Rees order.  The negative normalization is not merely bounded: its
ratio to the Wold volume converges exactly to
\(\zeta(1-c)^2/\zeta(1-2c)\), because Hellinger affinity is the midpoint Jensen ratio of the
finite Euler partition.

This three-phase law refines Kakutani singularity by recovering the orientation that ordinary
product-measure equivalence forgets.  More generally, all-orders flatness is a universal
superselection radical: every fixed finite-Rees operator preserves it in arbitrary dimension,
and every exact sewing to a nonzero finite-grade germ has norm beyond all inverse Rees powers.
The prime source is an arithmetic realization of that theorem, not an assumption built into it.
The same source has an independent depth--support superselection law: finite differential closure
cannot manufacture the complete valuation tower, and fusion escapes every finite prefix.

At a divisor, the only compatible escape is singular.  Differentiation drops one overlap grade,
positivity forces an inverse-square diagonal blow-up, and the exact M\"obius--Gamma overlap graph
uses that singularity to reconstruct the completed squarefree current and then the continued
\(\Xi\)-current.

The resulting universal lesson is precise.  Scalar reflection, finite jets, bounded
finite-order operators, and finite differential closures are shadows of a finer source geometry:
even action, odd tangent, depth--support filtration, and Rees asymptotic type.  Reflection either
preserves that typed tangent internally or pays modular drift in a complementary Weyl channel;
criticality is the unique zero-retuning locus.  At finite aperture the \(\Xi\) transport is
already four-unit nonexpansive.  Zero exclusion sits exactly at its singular boundary derivative.  A viable
source-internal RH architecture must therefore carry the pointed, divisor-graded valuation
tangent to that Ward boundary before scalarization erases its sign or grade.

\small
\begin{thebibliography}{99}

\bibitem{ArratiaBarbourTavare1997}
R.~Arratia, A.~D.~Barbour, and S.~Tavar\'e,
\emph{Random combinatorial structures and prime factorizations},
Notices Amer. Math. Soc. \textbf{44} (1997), 903--910.

\bibitem{ArratiaBarbourTavare2003}
R.~Arratia, A.~D.~Barbour, and S.~Tavar\'e,
\emph{Logarithmic Combinatorial Structures: A Probabilistic Approach},
EMS Monographs in Mathematics, European Mathematical Society, 2003.

\bibitem{CalegariVenkatesh}
F.~Calegari and A.~Venkatesh,
\emph{A torsion Jacquet--Langlands correspondence},
preprint, \S6.6.1,
\url{https://arxiv.org/abs/1212.3847}.

\bibitem{Connes1999}
A.~Connes,
\emph{Trace formula in noncommutative geometry and the zeros of the Riemann zeta function},
Selecta Math. (N.S.) \textbf{5} (1999), 29--106.

\bibitem{Edwards1974}
H.~M.~Edwards,
\emph{Riemann's Zeta Function},
Academic Press, 1974.

\bibitem{Hellinger1909}
E.~Hellinger,
\emph{Neue Begr\"undung der Theorie quadratischer Formen von unendlichvielen Ver\"anderlichen},
J. Reine Angew. Math. \textbf{136} (1909), 210--271.

\bibitem{Kakutani1948}
S.~Kakutani,
\emph{On equivalence of infinite product measures},
Ann. of Math. (2) \textbf{49} (1948), 214--224.

\bibitem{Jana2021}
S.~Jana,
\emph{Applications of analytic newvectors for \(\mathrm{GL}(n)\)},
Math. Ann. \textbf{380} (2021), 915--952.

\bibitem{KhelifScarpalezos2021}
A.~Khelif and D.~Scarpalezos,
\emph{Maximal closed ideals of the Colombeau algebra of generalized functions},
Monatsh. Math. \textbf{194} (2021), 119--150.

\bibitem{Mertens1874}
F.~Mertens,
\emph{Ein Beitrag zur analytischen Zahlentheorie},
J. Reine Angew. Math. \textbf{78} (1874), 46--62.

\bibitem{Ostrowski1920}
A.~Ostrowski,
\emph{\"Uber Dirichletsche Reihen und algebraische Differentialgleichungen},
Math. Z. \textbf{8} (1920), 241--298.

\bibitem{Pong2016}
W.-Y.~Pong,
\emph{Applications of differential algebra to algebraic independence of arithmetic functions},
Acta Arith. \textbf{172} (2016), 149--173.

\bibitem{Riemann1859}
B.~Riemann,
\emph{\"Uber die Anzahl der Primzahlen unter einer gegebenen Gr\"osse},
Monatsber. K\"onigl. Preuss. Akad. Wiss. Berlin (1859), 671--680.

\bibitem{ShapiroSparer1986}
H.~N.~Shapiro and G.~Sparer,
\emph{On algebraic independence of Dirichlet series},
Comm. Pure Appl. Math. \textbf{39} (1986), 695--745.

\bibitem{Titchmarsh1986}
E.~C.~Titchmarsh,
\emph{The Theory of the Riemann Zeta-Function}, 2nd ed.,
revised by D.~R.~Heath-Brown, Oxford University Press, 1986.

\end{thebibliography}

\end{document}
